\begin{helper}[Terminal-ladder interchangeability away from the PE pair]
Under the terminal parallel hypotheses and the condition (5.1) induction
context (lower assertions (i)--(vi) and current assertions (iii)--(vi)), let
$v_n,\gamma_n,\delta_n$ be the exact support and color fields of
the last PE step.  If $\alpha$ and $\beta$ are both missing on $T_n$ under
$\sigma_n$ and neither belongs to $\{\gamma_n,\delta_n\}$, then they are
$T_n$-interchangeable under $\sigma_n$, with isolated path components
included.
\end{helper}
\begin{proof}
Apply current assertion (iv) from \noderef{TechnicalInductionContext} to the
given terminal PE, the exact support and color identities, and the proper
stable coloring $\sigma_n$.  It says that the
$(\gamma_n,\delta_n)$-component through $v_n$ meets $V(T_n)$ only at
$v_n$.  Therefore \noderef{ParallelExtensionUndoStable} applies to the
component swap
\[
 \tau=\sigma_n/P_{v_n}(\gamma_n,\delta_n,\sigma_n)
\]
and makes $\tau$ stable on the preceding terminal ladder data
$(T_n,D_{n-1},\varphi_{n-1})$.

The PE certificate in \noderef{TashkinovSeriesStep} gives
$\gamma_n\ne\delta_n$, so
\noderef{TwoColorComponentSwapProper} transports properness from
$\sigma_n$ to $\tau$.  The definition \noderef{TwoColorComponentSwap}
fixes every uncolored value and sends colored values to colored values;
hence \noderef{IsOneEdgeUncolored} is also preserved.  It changes only
values $\gamma_n$ and $\delta_n$.  Since $\alpha$ and $\beta$ avoid those
two colors, an exact edge has color $\alpha$ or $\beta$ under $\tau$ if and
only if it has that same color under $\sigma_n$.  Thus the two
$(\alpha,\beta)$ incidence graphs agree edge for edge.  In particular their
connected components, isolated components, path-component predicates, and
the missingness of $\alpha,\beta$ on $T_n$ agree.

If $n>1$, \noderef{CertifiedETTRestriction} at rung $n-1$ gives the actual
positive lower-rung maximum ETT whose next ladder tree is $T_n$.  Lower
assertion (ii) in \noderef{TechnicalInductionContext}, applied to the stable
coloring $\tau$, gives $T_n$-interchangeability under $\tau$.  If $n=1$,
the interchangeability clause of \noderef{TechnicalTheoremBase} applied to
the original initial tuple gives the same conclusion.  Finally, equality of
the two incidence graphs transports every full component equality in the
definition \noderef{TreeInterchangeable} from $\tau$ back to $\sigma_n$.
This proves the assertion, including when one of the path components is an
isolated vertex.
\end{proof}
